\section{The number of final states in Hastings' hierarchical construction}
\label{sec:hastings_final_state_count}

% Hastings, arXiv:cond-mat/0508554v3, proj3new.tex, lines 722--730.

The hierarchical construction of \cite{Hastings2006Solving} ends with a family
of final states, and the source asserts that once each of them lies within
energy $\Delta E/2$ of the ground state there are at most $2n$ of them, $n$
being the number of ground states \cite[lines 722--730]{Hastings2006Solving}.
The source gives no argument; the following one rests only on the gap and on
Bessel's inequality \cite{gap:hastings05_final_state_count}.

\begin{lemma}[Bessel's inequality against a projection]
    \label{thm:hastings_final_state_bessel}
    \lean{TNLean.MPS.AreaLaw.LocalSolution.inner_toEuclideanCLM_eq_inner_self,
        TNLean.MPS.AreaLaw.LocalSolution.re_inner_toEuclideanCLM_eq_norm_sq,
        TNLean.MPS.AreaLaw.LocalSolution.sum_norm_sq_le_re_trace}
    \leanok
    Let $P$ be an orthogonal projection on $\mathbb C^d$ and let $\phi_k$,
    $k\in S$, be orthonormal. Then $\langle\phi_k,P\phi_k\rangle=\|P\phi_k\|^2$
    and
    \begin{align}
        \sum_{k\in S}\|P\phi_k\|^2\le\operatorname{tr}P.
    \end{align}
\end{lemma}

\begin{proof}
    \leanok
    Since $P$ is self-adjoint and $P^2=P$,
    $\langle\phi,P\phi\rangle=\langle\phi,P^2\phi\rangle=\|P\phi\|^2$. With
    $e_j$ the standard basis, $\|P\phi_k\|^2=\sum_j|\langle\phi_k,Pe_j\rangle|^2$.
    Exchanging the sums and applying Bessel's inequality to each $Pe_j$ gives
    $\sum_k\|P\phi_k\|^2\le\sum_j\|Pe_j\|^2=\sum_jP_{jj}=\operatorname{tr}P$.
\end{proof}

\begin{lemma}[Orthonormal states of low energy]
    \label{thm:hastings_final_state_count}
    \lean{TNLean.MPS.AreaLaw.LocalSolution.one_sub_mul_card_le_re_trace,
        TNLean.MPS.AreaLaw.LocalSolution.card_le_two_mul_re_trace}
    \leanok
    \uses{thm:hastings_final_state_bessel}
    Let $P$ be an orthogonal projection and $H$ an operator with
    $H-E_0\ge\Delta E\,(1-P)$ for some $\Delta E>0$, that is,
    $\Delta E\,(\|v\|^2-\langle v,Pv\rangle)\le\langle v,Hv\rangle-E_0\|v\|^2$
    for every $v$. If $\phi_k$, $k\in S$, are orthonormal with
    $\langle\phi_k,H\phi_k\rangle\le E_0+\theta\Delta E$, then
    \begin{align}
        (1-\theta)\,|S|\le\operatorname{tr}P.
    \end{align}
    In particular at most $2\operatorname{tr}P$ orthonormal states have energy
    at most $E_0+\Delta E/2$. With $P$ the projection onto the $n$ ground
    states this is the count of \cite[lines 722--730]{Hastings2006Solving},
    for final states whose energy under $H$ is within $\Delta E/2$ of the
    ground energy \cite{gap:hastings05_final_state_count}.
\end{lemma}

\begin{proof}
    \leanok
    \uses{thm:hastings_final_state_bessel}
    For each $k$ the gap gives
    $\Delta E\,(1-\|P\phi_k\|^2)\le\theta\Delta E$, so
    $\|P\phi_k\|^2\ge1-\theta$. Summing over $k$ and applying
    Lemma~\ref{thm:hastings_final_state_bessel} gives
    $(1-\theta)|S|\le\operatorname{tr}P$; take $\theta=1/2$ for the second
    statement.
\end{proof}
